% !TEX root =./ECE343_Practice_main.tex
%
% L14 practice set (Transients): turn-on, reflections at both ends, the bounce
% diagram, the steady state, and a pulse by superposition.

\fullwidth{\textbf{Learning Objective 3.7: } I can plot the transient response of a
load connected to a source through a transmission line.
\\[4pt]\normalfont\small Show your work. Write a \textbf{Documentation} line at the
end (\textbf{None} if you did not collaborate). All lines are lossless, and every
source switches on at $t = 0$.}

\begin{questions}

%%%%%%%%%%%%%%%%%%%%%%%%%%%%%%%%%%%%%%%%%
\question \textbf{Turn-on.} A $12\ \text{V}$ step source with internal resistance
$\myz{s} = 75\ \ohms$ drives a $300\ \m$ line with $\myz{0} = 50\ \ohms$ and
$v_p = 2\times10^{8}$ m/s. The line ends in $\myz{L} = 25\ \ohms$.
\begin{parts}

	\begin{minipage}[t]{0.58\linewidth}
		\part Find the one-way delay $T$.
		\begin{solution}[0.8in]
		\eq{ T = \dfrac{\ell}{v_p} = \dfrac{300}{2\times10^{8}} = 1.5\ \mu\text{s} }
		\end{solution}
	\end{minipage}\hfill%
	\begin{minipage}[t]{0.37\linewidth}
		\raggedleft \vspace{12pt}
		$T$ = \ansbox[1.3in]{$1.5\ \mu$s}
	\end{minipage}

	\begin{minipage}[t]{0.58\linewidth}
		\part Find the voltage $V_1^+$ launched onto the line at turn-on, and explain
		why the load does not affect it.
		\begin{solution}[1.2in]
		\eq{ V_1^+ = V_s\dfrac{\myz{0}}{\myz{s} + \myz{0}} = 12\times\dfrac{50}{125} = 4.8\ \text{V} }
		Just after turn-on only the incident wave exists, so the line looks like
		$\myz{0}$ to the source. The load has no effect until the step reaches it at
		$t = T$.
		\end{solution}
	\end{minipage}\hfill%
	\begin{minipage}[t]{0.37\linewidth}
		\raggedleft \vspace{12pt}
		$V_1^+$ = \ansbox[1.3in]{$4.8\ \text{V}$}
	\end{minipage}

	\begin{minipage}[t]{0.58\linewidth}
		\part Find $\Gamma_L$ and $\Gamma_s$.
		\begin{solution}[1in]
		\eq{
			\Gamma_L &= \dfrac{25-50}{25+50} = -\dfrac{1}{3} \\
			\Gamma_s &= \dfrac{75-50}{75+50} = 0.2
		}
		\end{solution}
	\end{minipage}\hfill%
	\begin{minipage}[t]{0.37\linewidth}
		\raggedleft \vspace{12pt}
		$\Gamma_L$ = \ansbox[1.3in]{$-0.333$}

		\vspace{4pt}
		$\Gamma_s$ = \ansbox[1.3in]{$0.2$}
	\end{minipage}

	\begin{minipage}[t]{0.58\linewidth}
		\part Find the final, steady-state voltage at the load.
		\begin{solution}[0.8in]
		\eq{ V_L(\infty) = V_s\dfrac{\myz{L}}{\myz{L} + \myz{s}} = 12\times\dfrac{25}{100} = 3.0\ \text{V} }
		\end{solution}
	\end{minipage}\hfill%
	\begin{minipage}[t]{0.37\linewidth}
		\raggedleft \vspace{12pt}
		$V_L(\infty)$ = \ansbox[1.3in]{$3.0\ \text{V}$}
	\end{minipage}

\end{parts}

\newpage
%%%%%%%%%%%%%%%%%%%%%%%%%%%%%%%%%%%%%%%%%
\question \textbf{The bounce diagram.} Continue with the line in Question 1.
\begin{parts}

	\part Draw the bounce diagram through $t = 5T$, labeling each wave with its
	amplitude.
	\begin{solution}[1.4in]
	Each wave is the previous one times the reflection coefficient at the end it
	reached:
	\eq{ 4.8,\quad -\tfrac{1}{3}(4.8) = -1.6,\quad 0.2(-1.6) = -0.32,\quad -\tfrac{1}{3}(-0.32) = 0.107,\quad 0.2(0.107) = 0.021\ \text{V} }
	Forward waves leave the source at $0$, $2T$, and $4T$; reflected waves leave
	the load at $T$ and $3T$.
	\end{solution}

	\begin{minipage}[t]{0.58\linewidth}
		\part Find the load voltage on each interval: $T < t < 3T$, $3T < t < 5T$, and
		$5T < t < 7T$.
		\begin{solution}[1.3in]
		\eq{
			V_L &= 4.8 - 1.6 = 3.20\ \text{V} \\
			V_L &= 3.20 - 0.32 + 0.107 = 2.99\ \text{V} \\
			V_L &= 2.99 + 0.021 - 0.007 = 3.00\ \text{V}
		}
		\end{solution}
	\end{minipage}\hfill%
	\begin{minipage}[t]{0.37\linewidth}
		\raggedleft \vspace{12pt}
		$V_L(2T)$ = \ansbox[1.3in]{$3.20\ \text{V}$}

		\vspace{4pt}
		$V_L(4T)$ = \ansbox[1.3in]{$2.99\ \text{V}$}

		\vspace{4pt}
		$V_L(6T)$ = \ansbox[1.3in]{$3.00\ \text{V}$}
	\end{minipage}

	\begin{minipage}[t]{0.58\linewidth}
		\part Find the voltage at the source end of the line on $0 < t < 2T$ and
		$2T < t < 4T$.
		\begin{solution}[1in]
		\eq{
			V(-\ell) &= 4.8\ \text{V} \\
			V(-\ell) &= 4.8 - 1.6 - 0.32 = 2.88\ \text{V}
		}
		\end{solution}
	\end{minipage}\hfill%
	\begin{minipage}[t]{0.37\linewidth}
		\raggedleft \vspace{12pt}
		$V(T)$ = \ansbox[1.3in]{$4.8\ \text{V}$}

		\vspace{4pt}
		$V(3T)$ = \ansbox[1.3in]{$2.88\ \text{V}$}
	\end{minipage}

\end{parts}

\newpage
%%%%%%%%%%%%%%%%%%%%%%%%%%%%%%%%%%%%%%%%%
\question \textbf{The steady state.} The load voltage after many round trips is
\eq{ V_L(\infty) = V_1^+\lp 1 + \Gamma_L\rp\sum_{n=0}^{\infty}\lp\Gamma_L\Gamma_s\rp^{n} }
\begin{parts}

	\part Explain why the series converges for any pair of passive resistive
	terminations other than both ends being shorts or opens.
	\begin{solution}[0.9in]
	For a passive resistor, $-1 \le \Gamma \le 1$, with $\vert\Gamma\vert = 1$ only
	for a short or an open. If either end is a resistor, $\vert\Gamma_L\Gamma_s\vert < 1$,
	and a geometric series with ratio less than one in magnitude converges.
	\end{solution}

	\part Sum the series and show that
	$V_L(\infty) = V_s\dfrac{\myz{L}}{\myz{L} + \myz{s}}$.
	\begin{solution}[2.2in]
	\eq{
		V_L(\infty) &= \dfrac{V_1^+\lp 1 + \Gamma_L\rp}{1 - \Gamma_L\Gamma_s} \\
		1 + \Gamma_L &= \dfrac{2\myz{L}}{\myz{L} + \myz{0}}, \qquad
		1 - \Gamma_L\Gamma_s = \dfrac{2\myz{0}\lp\myz{L} + \myz{s}\rp}{\lp\myz{L} + \myz{0}\rp\lp\myz{s} + \myz{0}\rp} \\
		V_L(\infty) &= V_s\dfrac{\myz{0}}{\myz{s} + \myz{0}}\cdot\dfrac{2\myz{L}}{\myz{L} + \myz{0}}\cdot
		\dfrac{\lp\myz{L} + \myz{0}\rp\lp\myz{s} + \myz{0}\rp}{2\myz{0}\lp\myz{L} + \myz{s}\rp}
		= V_s\dfrac{\myz{L}}{\myz{L} + \myz{s}}
	}
	\end{solution}

	\begin{minipage}[t]{0.58\linewidth}
		\part Check the result numerically with the values from Question 1.
		\begin{solution}[0.9in]
		\eq{ \dfrac{4.8\lp 1 - \tfrac{1}{3}\rp}{1 + \tfrac{1}{3}(0.2)} = \dfrac{3.2}{1.0667} = 3.0\ \text{V} }
		\end{solution}
	\end{minipage}\hfill%
	\begin{minipage}[t]{0.37\linewidth}
		\raggedleft \vspace{12pt}
		$V_L(\infty)$ = \ansbox[1.3in]{$3.0\ \text{V}$}
	\end{minipage}

\end{parts}

\newpage
%%%%%%%%%%%%%%%%%%%%%%%%%%%%%%%%%%%%%%%%%
\question \textbf{A pulse.} The worked example from the lesson ($V_s = 10\ \text{V}$,
$\myz{s} = 25\ \ohms$, $\myz{0} = 50\ \ohms$, $\myz{L} = 150\ \ohms$, $T = 2\ \mu$s)
is driven by a $10\ \text{V}$ pulse that lasts from $t = 0$ to $t = T$ instead of a
step.
\begin{parts}

	\part Write the pulse as the sum of two steps, and explain why the load voltage
	is the step response minus a delayed copy of itself.
	\begin{solution}[1in]
	The pulse is a $10\ \text{V}$ step at $t = 0$ plus a $-10\ \text{V}$ step at
	$t = T$. The line is linear, so its response is the sum of the responses to the
	two steps: $V_L(t) = s(t) - s(t - T)$, where $s(t)$ is the step response.
	\end{solution}

	\begin{minipage}[t]{0.58\linewidth}
		\part The step response at the load is $0$, then $10.00$, $8.33$, and
		$8.61\ \text{V}$ on successive intervals of $2T$ starting at $T$. Find the
		pulse response at the load on $T < t < 2T$, $2T < t < 3T$, and $3T < t < 4T$.
		\begin{solution}[1.2in]
		\eq{
			T < t < 2T&: \quad 10.00 - 0 = 10.0\ \text{V} \\
			2T < t < 3T&: \quad 10.00 - 10.00 = 0\ \text{V} \\
			3T < t < 4T&: \quad 8.33 - 10.00 = -1.67\ \text{V}
		}
		\end{solution}
	\end{minipage}\hfill%
	\begin{minipage}[t]{0.37\linewidth}
		\raggedleft \vspace{12pt}
		$V_L(1.5T)$ = \ansbox[1.3in]{$10.0\ \text{V}$}

		\vspace{4pt}
		$V_L(2.5T)$ = \ansbox[1.3in]{$0\ \text{V}$}

		\vspace{4pt}
		$V_L(3.5T)$ = \ansbox[1.3in]{$-1.67\ \text{V}$}
	\end{minipage}

	\part The load is a digital receiver. The pulse should produce one clean pulse
	at the load, but a second, inverted pulse arrives at $3T$. What single change to
	the circuit would remove it, and why?
	\begin{solution}[1in]
	Make the source resistance equal to $\myz{0} = 50\ \ohms$. The second pulse is the
	first reflection from the load, re-reflected by the source; with
	$\Gamma_s = 0$ the source absorbs it. (A $50\ \ohms$ load would also work,
	by removing the first reflection, but it would halve the received pulse.)
	\end{solution}

\end{parts}

\vspace{1.5em}
\noindent\textbf{Documentation:} \rule[-2pt]{0.6\linewidth}{0.4pt}

\end{questions}
